\section{Cuantitativas}
Las variables cuantitativas son aquellas que pueden ser medidas y expresadas numéricamente. Estas
variables permiten realizar operaciones matemáticas y estadísticas. A diferencia de la sección
anterior, en este caso podremos desarrollar una análisis más rico.

\subsection{Discretas}
La variable a analizar es \Variable{'absence\_days'}.

\TablaDeDatos
\begin{DatosTabla}{L{0.20\linewidth} C{0.14\linewidth} C{0.14\linewidth} C{0.14\linewidth} C{0.14\linewidth}}{\textbf{X (Días)} & \textbf{fi} & \textbf{fr} & \textbf{Fi} & \textbf{Fr}}

\TablaClave{0} & 36 & 32,43\% & 36 & 32,43\% \\
\TablaClave{1} & 31 & 27,93\% & 67 & 60,36\% \\
\TablaClave{2} & 27 & 24,32\% & 94 & 84,68\% \\
\TablaClave{3} & 17 & 15,32\% & 111 & 100,00\% \\
\textbf{Total} & \textbf{111} & \textbf{100,00\%} & \textbf{111} & \textbf{100,00\%} \\

\end{DatosTabla}


\AnalisisDescriptivo
\begin{DatosTabla}{L{0.24\linewidth} C{0.14\linewidth} L{0.44\linewidth}}{\textbf{Medida} & \textbf{Valor} & \textbf{Interpretación}}[Analisis descriptivo de absence\_days]

Muestra (N) & 111 & Total de estudiantes analizados. \\ \hline
Media & 1,22 & Promedio de días de ausencia por alumno. \\ \hline
Mediana & 1 & El 50\% de los alumnos tiene 1 falta o menos. \\ \hline
Moda & 0 & El valor más frecuente. \\ \hline
Desviación estándar & 1,1 & Dispersión de los datos respecto a la media. \\ \hline
Mínimo & 0 & Menor cantidad de faltas registradas. \\ \hline
Máximo & 3 & Mayor cantidad de faltas registradas. \\ \hline
Rango & 3 & Amplitud total de los datos. \\


\end{DatosTabla}



\TituloGrafico{ausencias}
% =====================================================================
%  figures/ausencias-bastones.tex
%  Grafico de BASTONES para la variable discreta 'absence_days'.
%  (Reemplaza a img/discreta.png; correccion de la profesora:
%   para variables discretas corresponde grafico de bastones o de caja)
% =====================================================================
\begin{figure}[H]
\centering
\begin{tikzpicture}
\begin{axis}[
  ybar, bar width=22pt,
  title={Frecuencia de días de ausencia (\Variable{absence\_days})},
  xlabel={Días de ausencia},
  ylabel={Cantidad de estudiantes (fi)},
  ymin=0, ymax=42,
  symbolic x coords={0, 1, 2, 3},
  xtick=data,
  enlarge x limits=0.4,
  nodes near coords,
]
\addplot+[fill=azul, draw=azul] coordinates {(0,36) (1,31) (2,27) (3,17)};
\end{axis}
\end{tikzpicture}
\caption{Gráfico de bastones para la variable \Variable{absence\_days}.}
\label{fig:grafico_absence_days}
\end{figure}


\Interpretacion
Este análisis evalúa el ausentismo, registrando una asistencia general bastante regular y un alto
compromiso con las clases. El promedio de inasistencias es sumamente bajo (1,22 días). Que la moda
sea 0 indica que el escenario más común en el aula es la asistencia perfecta. La dispersión es
mínima, ya que la brecha total entre el que más faltó y el que menos lo hizo es de apenas 3 días.

%---------------------------------------------------------------------------------------------------

\subsection{Continuas}
La variable a analizar es \Variable{'math\_score'}.

\AnalisisDescriptivo
\begin{DatosTabla}{L{0.24\linewidth} C{0.14\linewidth} L{0.44\linewidth}}{\textbf{Medida} & \textbf{Valor} & \textbf{Interpretación}}[Analisis descriptivo de math\_score]

Moda & 85 & Para la variable \Variable{math\_score}, el valor que más se repite es 85. \\ \hline
Media & 85,09 & En promedio los alumnos obtuvieron 85,09 como puntaje de matemática. \\ \hline
Mediana & 82 & El 50\% de los estudiantes sacó 82 o menos. \\ \hline
Primer Cuartil & 80 & El 25\% de los estudiantes sacó 80 o menos. \\ \hline
Tercer Cuartil & 85 & El 75\% de los estudiantes sacó 85 o menos. \\ \hline
Varianza & 10,62 & Los puntajes presentan una dispersión baja, ya que la mayoría de los estudiantes obtuvo resultados cercanos al promedio de 85,09 puntos. \\ \hline
Desvío Estándar & 3,26 & Los puntajes suelen diferir de la media en aproximadamente 3,26 puntos. \\ \hline
Rango Intercuartílico & 5,00 & En el 50\% central hay una diferencia de 5 puntos. \\ \hline
Coeficiente de Variación & 3,83\% & Como es menor al 20\%, se puede decir que los datos son homogéneos y, por lo tanto, la media es representativa de este conjunto de datos. \\


\end{DatosTabla}



\Interpretacion
La moda y la mediana son cercanas a la media, lo que indica que la mayoría de los estudiantes obtuvo
resultados similares. El rango intercuartílico es de 5 puntos, lo que sugiere que el 50\% central de
los puntajes se encuentra dentro de un rango estrecho. El coeficiente de variación del 3,83\%
confirma la homogeneidad de los datos, lo que hace que la media sea representativa del conjunto de
datos.

\subsubsection{Análisis por intervalos}
Los resultados previamente mostrados surgen de un análisis con los datos sin agrupar. Nos propusimos
plantear una nueva perspectiva, agrupando los puntajes en intervalos y comparar los resultados de
ambos análisis.

\begin{DatosTabla}{C{0.16\linewidth} C{0.12\linewidth} C{0.12\linewidth} C{0.12\linewidth} C{0.12\linewidth} C{0.12\linewidth}}{\textbf{Intervalo} & \textbf{fi} & \textbf{fr} & \textbf{Fi} & \textbf{Fr} & \textbf{MC}}[Frecuencias de math\_score por intervalos]

74 - 78 & 1 & 0,90\% & 1 & 0,90\% & 76 \\ \hline
78 - 82 & 10 & 9,01\% & 11 & 9,91\% & 80 \\ \hline
82 - 86 & 64 & 57,66\% & 75 & 67,57\% & 84 \\ \hline
86 - 90 & 21 & 18,92\% & 96 & 86,49\% & 88 \\ \hline
90 - 94 & 15 & 13,51\% & 111 & 100,00\% & 92 \\ \hline
\textbf{Total} & \textbf{111} & \textbf{100,00\%} & \textbf{111} & \textbf{100,00\%} & -- \\


\end{DatosTabla}


\PalabraClave{Media}: 85,34 -
\PalabraClave{Varianza}: 10,06 -
\PalabraClave{Desvío estándar}: 3,23 -
\PalabraClave{CV}: 3,78\%

\TituloGrafico{Histograma con intervalos}
% =====================================================================
%  figures/intervalos-histograma.tex
%  Histograma de 'math_score' agrupada en intervalos.
%  (Reemplaza a img/intervalos.png)
% =====================================================================
\begin{figure}[H]
\centering
\begin{tikzpicture}
\begin{axis}[
  title={Distribución de \Variable{math\_score} por intervalos},
  xlabel={Intervalo de puntaje},
  ylabel={Frecuencia (fi)},
  ymin=0, ymax=75,
  xmin=74, xmax=94,
  xtick={76,80,84,88,92},
  xticklabels={74-78, 78-82, 82-86, 86-90, 90-94},
  enlarge x limits=false,
]
\addplot[ybar interval, fill=azul, draw=azul] coordinates {
  (74,1) (78,10) (82,64) (86,21) (90,15) (94,0)
};
% Etiquetas de frecuencia centradas sobre cada barra (barras contiguas = histograma)
\node[anchor=south, text=azul] at (axis cs:76,1) {1};
\node[anchor=south, text=azul] at (axis cs:80,10) {10};
\node[anchor=south, text=azul] at (axis cs:84,64) {64};
\node[anchor=south, text=azul] at (axis cs:88,21) {21};
\node[anchor=south, text=azul] at (axis cs:92,15) {15};
\end{axis}
\end{tikzpicture}
\caption{Histograma de la variable \Variable{math\_score} agrupada en intervalos.}
\label{fig:grafico_math_score_intervalos}
\end{figure}


\subsubsection{Tabla comparativa}
\begin{DatosTabla}{L{0.34\linewidth} C{0.30\linewidth} C{0.32\linewidth}}{\textbf{Medida} & \textbf{Análisis con datos sin agrupar} & \textbf{Análisis por intervalos}}[Comparacion: datos sin agrupar vs por intervalos]

Media & 85,09 & 85,34 \\ \hline
Varianza & 10,62 & 10,06 \\ \hline
Desvío estándar & 3,26 & 3,23 \\ \hline
CV & 3,83\% & 3,78\% \\


\end{DatosTabla}


\subsubsection{Conclusión}
Al comparar ambos análisis, se observa que las medidas obtenidas a partir de los datos agrupados son
muy similares a las calculadas con los datos originales, aunque presentan una ligera pérdida de
precisión debido al uso de intervalos y marcas de clase. Sin embargo, la agrupación facilita la
lectura e interpretación de la distribución de los datos, permitiendo identificar rápidamente su
comportamiento general sin necesidad de analizar cada observación individual.
